\section{Introduction}

Theorem D.1 of OpenAI's 25 October 2026 manuscript~\cite{OpenAI26}
asserts a deterministic, polynomial-space contraction algorithm for
Boolean tensor networks with linearly many edges and maximum degree
$\lceil s^{9/10}\rceil$. We denote this assertion by \HTC, and state its
full interface in Section~\ref{sec:interface}. This paper studies its
consequences. The degree bound, the allowed coefficient ring, the cost
of evaluating implicit tensors, and exact integer readback are part of
every application.

Throughout, $\log_2$ is the binary logarithm, and
$g(s)=s/\sqrt{\log_2(s+16)}$. Running times count bit operations, including
polynomial factors in the description length. Polynomial space means
polynomially many bits. A circuit wire is an occurrence in a gate's
input list; arbitrary fan-out is permitted. Circuit depth and fan-in
are unrestricted unless a statement explicitly bounds them.

\subsection{Statements of the consequences}

\begin{theorem}[Threshold acceptance probability, conditional on \HTC]
\label{thm:capp}
Assume \HTC. Fix $a\ge1$. Let $C$ be a single-output Boolean threshold
circuit with $n$ declared inputs, $w$ wires, and integer weights and
thresholds in a complete description of length $L\le S^a$, where
$S=n+w+2$. For every
\[
 2^{-S^{9/10}/\log_2(S+16)}\le\eps<1,
\]
there is a deterministic algorithm producing a rational $\widehat p$
with
\[
 \left|\widehat p-2^{-n}|\{x:C(x)=1\}|\right|\le\eps
\]
in $\poly_a(S+L)2^{O_a(g(S))}$ time and $\poly_a(S+L)$ space.
In particular this holds for $\eps\ge S^{-a}$. The algorithm can use the
fixed lower endpoint as its target error, so no arbitrary-precision real
encoding of $\eps$ is required.
\end{theorem}

The proof is in Sections~\ref{sec:protocol}--\ref{sec:unified}.
Theorem~\ref{thm:unified} gives the precise per-gate version. It allows
$\THR$, $\ETHR$, and weighted $\MOD_m$ gates, as well as small gates
with polynomial-time evaluators. Each large gate is assigned either an
exact sliced gadget, an exact supplied weight-factorization gadget, or
a hashed gadget. Under the explicit wire budget~\eqref{eq:budget},
the same time and space bounds hold, with error at most the sum of the
hashed gates' error budgets. When that set is empty, the answer is the
exact number of accepting assignments. In that exact-only case
the description length $L$ is arbitrary, with polynomial cost in $S+L$.
In particular, exact counting is available when the wire-weighted sum
of large gates' parameter bit lengths is $O(S^{19/10})$.

\begin{theorem}[Real-weight wire lower bounds, conditional on \HTC]
\label{thm:lower}
Assume \HTC. There are a constant $c>0$ and a language
$L_{\mathrm{ae}}\in\ENP$ such that, for every sufficiently large $N$,
every real-weight threshold circuit computing
$(L_{\mathrm{ae}})_N$ has more than $cN\sqrt{\log_2N}$ wires.
Consequently there is also such a language with this lower bound at
infinitely many lengths. For every fixed $\gamma\in(0,1/2)$ there are
$c_\gamma>0$ and $L_\gamma\in\ENP$ such that, for every sufficiently
large $N$, every real-weight threshold circuit with at most
$c_\gamma N\sqrt{\log_2N}$ wires agrees with $(L_\gamma)_N$ on at most
a $1/2+\gamma$ fraction of inputs. The circuits have arbitrary depth
and fan-in. The language in the last assertion may depend on $\gamma$.
\end{theorem}

Section~\ref{sec:lower} proves these statements by substituting the
acceptance-probability algorithm into the proofs of known transfer
constructions. We also give the separate infinitely-often construction,
since it makes the two-stage verification cost explicit.

\begin{theorem}[Undirected sparse connectivity, conditional on \HTC]
\label{thm:undirected}
Assume \HTC. Fix $c>0$. On simple undirected graphs with $n$ vertices
and $m\le cn$ edges, Hamiltonian cycle, Hamiltonian path (with specified
or unspecified endpoints), longest simple path, Hamiltonian TSP,
Steiner tree, feedback vertex set, connected vertex cover, and connected
dominating set admit deterministic decision, optimization, and witness
recovery in
\begin{equation}\label{eq:graph-time}
 \poly_c(n+L)2^{O_c(g(n))}\quad\hbox{time and}\quad\poly_c(n+L)
 \quad\hbox{space}.
\end{equation}
The relevant vertex and edge costs may be nonnegative integers of at
most $\lceil n^{9/10}\rceil$ bits. One can also exactly count undirected
Hamiltonian cycles as edge sets, Hamiltonian paths with specified
endpoints, simple paths with specified endpoints and length, and
Steiner tree subgraphs $(U,F)$. For connected covers, connected
dominating sets, and feedback vertex sets, the proof uses positive
weighted counts to decide existence and optimize.
\end{theorem}

\begin{theorem}[Directed sparse connectivity, conditional on \HTC]
\label{thm:directed}
Assume \HTC. Fix $c>0$. On digraphs without loops or parallel arcs and
with $m\le cn$, the bounds~\eqref{eq:graph-time} hold for exact directed
Hamiltonian-cycle counting, exact counting of directed paths with
specified endpoints and length, longest directed simple path,
nonnegative-cost directed Steiner tree, and minimum strongly connected
spanning subgraph. They also hold for a minimum strongly connected
subgraph containing specified terminals. Directed Steiner arborescences
with a fixed root, direction, and terminal set can be counted exactly.
Costs have at most $\lceil n^{9/10}\rceil$ bits. Optimization includes
recovery of an optimal subgraph.
\end{theorem}

\begin{theorem}[Routing on a sparse base graph, conditional on \HTC]
\label{thm:routing}
Assume \HTC. Fix $c>0$. Given a simple undirected base graph with
$m\le cn$, nonnegative integer edge costs of at most
$\lceil n^{9/10}\rceil$ bits, a terminal set $T$, and a required edge
set $R$, a minimum-cost closed walk visiting $T$ and traversing $R$
can be found, or infeasibility detected, within~\eqref{eq:graph-time}.
This includes graphical TSP, Steiner TSP, and undirected Rural Postman.
On a connected base graph, $T=V$ and $R=\varnothing$ give TSP on its
shortest-path metric, with an optimal metric tour recovered for $n\ge3$.
\end{theorem}

\subsection{Comparison and scope}

Table~\ref{tab:comparison} compares exponent scales. In its previous-best
column, $2^{\Theta(n)}$ denotes the single-exponential scale of the
available generic bounds; different problems and graph subclasses have
different constant bases. The cited algorithms specify representative
bounds within that scale. The comparison concerns unrestricted depth in
the circuit rows and a fixed edge-to-vertex ratio in the graph rows.
\begin{table}[ht]
\centering\small
\renewcommand{\arraystretch}{1.18}
\begin{tabular}{@{}P{.29\linewidth}P{.32\linewidth}P{.32\linewidth}@{}}
\toprule
Problem & Previous best (exponent scale) & Under \HTC\\
\midrule
Arbitrary-depth THR CAPP, $w=O(n)$ & $2^{\Theta(n)}$ generic enumeration & $2^{O(g(n))}$, error down to $2^{-\Theta(n^{.9}/\log n)}$\\
Mixed-gate exact counting under~\eqref{eq:budget} & $2^n\poly(L)$ enumeration & $\poly(S+L)2^{O(g(S))}$ for arbitrary $L$ if no gate is hashed\\
Unrestricted-depth real THR wires for explicit $\ENP$ languages & Linear scale; depth-restricted bounds differ & $>cN\sqrt{\log N}$ a.e.; fixed-advantage average case\\
Sparse Hamiltonicity & Single exponential; e.g. randomized $O^*(1.657^n)$~\cite{Bjorklund14} & Deterministic $2^{O(g(n))}$; exact \#HC\\
Sparse paths, Steiner, FVS, connected cover/domination & $2^{\Theta(n)}$ scale; connectivity frameworks~\cite{CutCount11,BCKN15} & Deterministic $2^{O(g(n))}$, optimization and witnesses\\
Sparse directed connectivity and sparse-base routing & $2^{\Theta(n)}$ scale (enumerate selected arcs or edge multiplicities) & Deterministic $2^{O(g(n))}$\\
\bottomrule
\end{tabular}
\caption{Polynomial factors in the input length are suppressed.
Algorithmic comparisons are deterministic unless randomization is marked. The
conditional bounds retain polynomial space. The exact-counting rows use
the object conventions in the theorem statements.}
\label{tab:comparison}
\end{table}

The polynomial description condition in Theorem~\ref{thm:capp} is a bit
length condition. For lower bounds, each real-weight threshold gate has
an equivalent polynomial-bit integer representation, with no change
in wires. This is an existence assertion, sufficient for guessing
circuits; it does not prescribe how to read arbitrary real-number
oracles. For an algorithm on a larger wire budget, the displayed bound
is in $S$, and may exceed exhaustive enumeration over the $n$ inputs.

\subsection{Related work}

OpenAI's manuscript itself derives a circuit lower bound from the
fixed-density case of its Theorem~1.1~\cite[Corollary~1.4]{OpenAI26}:
there are $L\in\ENP$ and $c>0$ such that circuits of fan-in-two Boolean
gates computing $L$ need more than $cn\sqrt{\log n}$ noninput gates for
infinitely many $n$. Its proof combines Williams' witness-circuit
argument with a local reduction to 3-CNF. Theorem~\ref{thm:lower} uses the
stronger input \HTC\ and a different measure. Every fan-in-two Boolean
gate has a threshold implementation with at most five wires, so
Theorem~\ref{thm:lower} implies a gate lower bound of the same order for
fan-in-two circuits, at every sufficiently large length and in the
constant-advantage form. In the other direction, the generic simulation
of Section~\ref{sec:lower} replaces a fan-in-$r$ threshold gate by
$O(r^2\log(r+2))$ fan-in-two gates, so Corollary~1.4 does not yield the
wire bound of Theorem~\ref{thm:lower} by gate simulation.

Williams' algorithmic method~\cite{Williams13} connects savings in
circuit analysis to circuit lower bounds. Chen and Williams~\cite{CW19}
use PCPs of proximity to reduce the composition overhead in such
connections. Bathie and Williams~\cite[Section 3.2--3.3]{BW24} give a
two-stage PCP/ECC/PCPP construction whose circuit queries are small
compositions of proof-bit circuits. The infinitely-often argument here
is a routine instantiation of that proof framework in the wire measure.
Chen, Lyu, and Williams~\cite{CLW20} give the refuter for a hierarchy
with sublinear nondeterminism and an XOR amplification based on linear
sums. Their proof mechanisms supply the almost-everywhere and
constant-advantage arguments below; the polynomial parameters and the
per-term wire counts are recalculated explicitly.

Depth-restricted threshold bounds have different quantitative behavior.
Impagliazzo, Paturi, and Saks~\cite{IPS97} prove that parity requires
$n^{1+c\theta^{-d}}$ wires at depth $d$, for absolute constants $c>0$
and $\theta\le3$. Kane and Williams~\cite{KW16} prove, for a fixed
advantage, depth-two lower bounds of order
$n^{5/2}/\log^{7/2}n$ wires and $n^{3/2}/\log^3n$ gates for an explicit
function; their depth-three result has a majority output gate.
Chen, Tal, and Wang's 2026 manuscript~\cite[Theorems 1.1--1.2]{CTW26}
states $n^{2.5-\eps}$ gate lower bounds in $\ENP$ for depth-two linear
threshold circuits, via a CAPP algorithm for an XOR of two such
circuits. These gate and wire measures and their depth restrictions
are retained in the comparison. Li and Yang~\cite{LiYang22} prove a
$3.1n-o(n)$ gate lower bound over the full binary Boolean basis. A
binary Boolean gate has a constant-size threshold implementation; the
reverse simulation of an arbitrary-weight threshold gate has a larger
bit cost, as quantified in Section~\ref{sec:lower}.

The comparison protocol used here follows Nisan's
integer-comparison construction~\cite{Nisan93}. Viola~\cite{Viola15}
improves the communication for addition and comparison. We use an
explicit half-splitting and prime-fingerprinting implementation to keep
seed length, local evaluation, and tensor ports in the same ledger.
The conversion from a deterministic transcript to local consistency
tensors and the union bound over the earliest erroneous gate are standard constructions.

For graph connectivity, Cut\&Count~\cite{CutCount11} combines consistent
cuts with isolation to filter disconnected candidates modulo two.
Bodlaender, Cygan, Kratsch, and Nederlof~\cite[Section 4]{BCKN15}
use determinants and the matrix-tree theorem for deterministic
connectivity counting parameterized by treewidth. Bj\"orklund's
Hamiltonicity algorithm~\cite{Bjorklund14} uses a different determinant
summation and cancellation. Our graph proofs instantiate the
matrix-tree approach using an explicit representation of the sparse
permutation sign by linearly many Boolean wires. The resulting scale is
$n+m$; no treewidth parameter is substituted for the number of vertices.
The $0/1/2$ multiplicity reduction for graphical TSP also appears in
the Cut\&Count treatment~\cite[Section 4.10]{CutCount11}.
